\chapter{Nissen Fundoplication}

\section*{What Is This Surgery?}
Nissen fundoplication is a definitive surgical procedure primarily designed to treat severe and chronic gastroesophageal reflux disease (GERD) by reinforcing the lower esophageal sphincter (LES) mechanism. This operation involves mobilizing the upper portion of the stomach, known as the fundus, and wrapping it completely (360 degrees) around the distal esophagus. The primary goal is to create a new, competent one-way valve at the gastroesophageal junction, thereby preventing the retrograde flow of gastric contents into the esophagus. It is typically performed when medical management with proton pump inhibitors (PPIs) fails to control symptoms, or when patients experience complications of GERD such as esophagitis, stricture, or Barrett's esophagus, or simply prefer a surgical solution to lifelong medication. This procedure aims to restore the anatomical integrity of the anti-reflux barrier and improve the patient's quality of life.

\section*{Alternative Names}
\begin{itemize}
  \item Laparoscopic Nissen Fundoplication (LNF)
  \item Total Fundoplication
  \item Anti-reflux Surgery
  \item Fundoplication
  \item GERD Surgery
  \item Nissen Wrap
\end{itemize}

\section*{Relevant Surgical Anatomy}
A thorough understanding of the anatomy of the gastroesophageal junction and surrounding structures is paramount for a successful Nissen fundoplication. The **distal esophagus** is approximately 2-4 cm in length below the diaphragm, traversing the **esophageal hiatus** of the diaphragm. It is anchored to the diaphragm by the **phrenoesophageal ligament**, a fibroelastic membrane that is often attenuated or disrupted in patients with hiatal hernias and GERD. The **diaphragmatic crura**, specifically the right crus, form the muscular sling around the esophagus at the hiatus, contributing to the extrinsic component of the lower esophageal sphincter. Repair of the often-widened hiatus by approximating the crura posteriorly is a critical step in restoring the anti-reflux barrier.

The **gastric fundus**, the uppermost part of the stomach, is the key component used to create the wrap. It lies superior and lateral to the gastroesophageal junction and is typically mobile, especially after division of the **short gastric vessels** that connect it to the splenic hilum. These vessels, usually 5-7 in number, originate from the splenic artery and vein. The **vagus nerves** (anterior and posterior trunks) descend along the esophagus and are intimately associated with its wall at the gastroesophageal junction. The anterior vagal trunk typically lies on the anterior surface of the esophagus, while the posterior vagal trunk is situated more posteriorly. Care must be taken to preserve these nerves, as injury can lead to gastric emptying disorders or diarrhea. The arterial supply to the distal esophagus and gastroesophageal junction is primarily from esophageal branches of the left gastric artery and inferior phrenic arteries, while venous drainage is via esophageal veins to the azygos/hemiazygos system and gastric veins to the portal system. Lymphatic drainage follows these vascular pathways to periesophageal and paracardial lymph nodes.

\section{Surgical Principle \& Rationale}
The fundamental principle of Nissen fundoplication is to restore the natural barrier function of the gastroesophageal junction, which is compromised in patients with GERD, often due to a hypotensive lower esophageal sphincter (LES) and/or a hiatal hernia. This is achieved by constructing a new, high-pressure zone that acts as a one-way valve, allowing food and liquids to pass into the stomach but preventing reflux of gastric acid and bile. The rationale is to mechanically augment the LES, thereby alleviating symptoms and allowing the esophageal mucosa to heal.

The procedure involves several key technical components:
\begin{itemize}
  \item \textbf{Esophageal Mobilization:} The distal esophagus is carefully dissected from its surrounding attachments within the mediastinum to ensure sufficient intra-abdominal length (typically 2-3 cm) for a tension-free wrap. This prevents the wrap from migrating back into the chest.
  \item \textbf{Hiatal Hernia Repair:} If a hiatal hernia is present, the enlarged diaphragmatic hiatus is meticulously closed with sutures, often reinforced with mesh in larger defects, to prevent upward migration of the stomach and maintain the integrity of the gastroesophageal junction. This step is crucial for long-term success.
  \item \textbf{Fundic Wrap Creation:} The gastric fundus is mobilized by dividing the short gastric vessels, allowing it to be brought behind the esophagus. It is then wrapped 360 degrees around the distal esophagus and secured to itself with sutures, typically over a large bougie (e.g., 54-60 Fr) placed in the esophagus. This creates a "floppy" wrap that allows for normal swallowing while effectively resisting reflux. The resulting wrap increases the pressure in the lower esophageal sphincter region, preventing gastric acid and bile from refluxing into the esophagus, thereby alleviating symptoms and allowing esophageal healing.
\end{itemize}

\section*{History \& Surgical Evolution}
The concept of surgically treating gastroesophageal reflux dates back to the early 20th century, but the modern era of fundoplication began with Rudolf Nissen. In 1955, while operating on a patient with a bleeding esophageal varices, Nissen inadvertently wrapped the gastric fundus around the esophagus to control the bleeding. He observed that this patient subsequently experienced complete resolution of severe heartburn symptoms. Nissen formally described his technique, the 360-degree fundoplication, in 1956, initially performing it via an open abdominal incision. For decades, the open Nissen fundoplication remained the gold standard for medically refractory GERD.

The next major revolution occurred in 1991 when Bernard Dallemagne and his team in Li\`{e}ge, Belgium, performed the first laparoscopic Nissen fundoplication. This minimally invasive approach significantly reduced postoperative pain, hospital stay, and recovery time, leading to a dramatic increase in the procedure's popularity and accessibility. Subsequent refinements have focused on optimizing wrap length and tension, leading to variations like the "floppy Nissen" (emphasizing a loose wrap over a large bougie) and the development of partial fundoplications (e.g., Toupet, Dor) for specific patient populations, particularly those with esophageal motility disorders or severe dysphagia concerns. The introduction of robotic platforms has further refined the laparoscopic approach, offering enhanced dexterity and 3D visualization.

\section*{Clinical Indications}
Nissen fundoplication is primarily indicated for patients suffering from chronic and severe gastroesophageal reflux disease (GERD) who meet specific criteria after a thorough diagnostic workup.
\begin{itemize}
  \item Persistent or recurrent GERD symptoms despite optimal and prolonged medical management with proton pump inhibitors (PPIs) for at least 8-12 weeks.
  \item Presence of GERD-related complications such as severe erosive esophagitis (Los Angeles grade C or D), esophageal stricture requiring repeated dilation, recurrent aspiration pneumonia, or Barrett's esophagus with low-grade dysplasia.
  \item Atypical or extra-esophageal manifestations of GERD, including chronic cough, asthma, hoarseness, laryngitis, or non-cardiac chest pain, where GERD is definitively proven to be the underlying cause by objective testing (e.g., 24-hour pH monitoring).
  \item Patients with a significant hiatal hernia (Type II, III, or IV) that contributes to reflux and is not adequately managed by medical therapy.
  \item Young patients who prefer to avoid lifelong medication, or those who experience intolerable side effects from PPIs, after a comprehensive discussion of risks and benefits.
  \item Predominant symptoms of regurgitation that are often less responsive to PPIs, especially when associated with a large hiatal hernia.
  \item Documented pathological acid reflux on 24-hour pH monitoring or pH-impedance testing, correlating with symptoms.
\end{itemize}

\section*{Clinical Positioning}
Nissen fundoplication is considered the primary surgical treatment for gastroesophageal reflux disease. In the broader context of GERD management, it is typically a secondary intervention, meaning it is usually pursued after a trial of conservative measures (lifestyle modifications) and medical therapy (e.g., proton pump inhibitors) has proven insufficient, undesirable, or has led to complications.

However, for certain patient profiles, it can be considered a primary treatment option. For instance, young, otherwise healthy individuals with documented severe GERD, a large hiatal hernia, or significant extra-esophageal symptoms, who wish to avoid lifelong medication, may opt for surgery as a first-line definitive treatment. It is also a primary choice for patients with complications like recurrent strictures or severe esophagitis that are not controlled by medical therapy, or those with significant regurgitation that is poorly responsive to PPIs. In essence, while medical therapy is generally the initial approach, Nissen fundoplication serves as the primary surgical solution when the disease warrants operative intervention, offering a durable anatomical correction.

\section{Surgical Technique \& Procedure}
Nissen fundoplication is most commonly performed using a minimally invasive laparoscopic approach, though an open technique may be necessary in select complex cases or for patients with extensive prior abdominal surgery. The procedure typically follows these steps:

\begin{enumerate}
  \item \textbf{Anesthesia and Positioning:} The patient is placed under general anesthesia. They are positioned supine with legs abducted (French position) or together, often in reverse Trendelenburg to allow gravity to retract abdominal contents inferiorly, improving exposure of the upper abdomen.
  \item \textbf{Access and Insufflation:} The abdomen is accessed, typically via an umbilical incision using a Veress needle or open Hasson technique, and insufflated with carbon dioxide to create a pneumoperitoneum (12-15 mmHg), providing working space. Multiple small incisions (usually 4-5, 5-12 mm) are then made for trocar placement to introduce laparoscopic instruments.
  \item \textbf{Liver Retraction:} The left lobe of the liver is gently retracted anteriorly and superiorly using a specialized retractor (e.g., Nathanson liver retractor) to expose the gastroesophageal junction and diaphragmatic hiatus.
  \item \textbf{Hiatal Dissection and Crural Repair:} The phrenoesophageal ligament is divided, and the distal esophagus is carefully mobilized from the surrounding mediastinal structures (anteriorly, posteriorly, and laterally) to achieve adequate intra-abdominal esophageal length (2-3 cm). Any existing hiatal hernia sac is reduced and excised. The diaphragmatic crura are then approximated posteriorly with non-absorbable sutures to narrow the hiatus, preventing future herniation and providing a stable platform for the wrap.
  \item \textbf{Fundic Mobilization:} The short gastric vessels, which connect the gastric fundus to the spleen, are often divided using energy devices (e.g., harmonic scalpel, LigaSure) to ensure adequate mobilization of the fundus for a tension-free wrap. This allows the fundus to be brought behind the esophagus without undue tension.
  \item \textbf{Bougie Placement and Wrap Creation:} A large esophageal bougie (e.g., 54-60 Fr) is inserted into the esophagus and stomach by the anesthesiologist to serve as a stent, preventing the wrap from being too tight and causing severe dysphagia. The mobilized gastric fundus is then passed behind the esophagus, creating a 360-degree wrap. The anterior wall of the fundus is then sutured to the posterior wall of the fundus, incorporating a small portion of the esophageal wall with each stitch to secure the wrap in place and prevent slippage. Typically, 2-3 non-absorbable sutures are placed.
  \item \textbf{Wrap Assessment and Closure:} The wrap is inspected for proper tension and length (typically 2-3 cm). The bougie is removed. Hemostasis is confirmed, and instruments are withdrawn. The port sites are closed in layers, with fascial closure for larger ports (typically >10 mm) to prevent incisional hernias, and skin closure using sutures or staples. No drains are typically placed unless there is a specific indication.
\end{enumerate}

\section*{Preoperative Workup \& Preparation}
Thorough preoperative preparation is crucial to ensure patient safety and optimize surgical outcomes for Nissen fundoplication. This comprehensive evaluation aims to confirm the diagnosis of GERD, rule out contraindications, and assess the patient's overall fitness for surgery.

\begin{itemize}
  \item \textbf{Diagnostic Workup:}
    \begin{itemize}
      \item \textbf{Upper Endoscopy (EGD):} To assess the esophageal and gastric mucosa, identify esophagitis, strictures, Barrett's esophagus, or other pathologies, and rule out malignancy. Biopsies are taken as indicated.
      \item \textbf{Esophageal Manometry:} To evaluate esophageal motility and rule out severe dysmotility disorders (e.g., achalasia, diffuse esophageal spasm, scleroderma) that would contraindicate a 360-degree Nissen fundoplication. It also helps identify a hypotensive LES.
      \item \textbf{24-hour pH Monitoring (or pH-impedance testing):} To objectively confirm the presence and severity of pathological acid reflux, especially in patients with atypical symptoms or equivocal endoscopic findings. This is critical for surgical planning.
      \item \textbf{Barium Swallow Study (Esophagram):} May be performed to assess esophageal anatomy, identify hiatal hernia size and type, evaluate for strictures, or assess for motility issues.
    \end{itemize}
  \item \textbf{Medical Clearance:} Comprehensive medical evaluation, including cardiac and pulmonary clearance, is obtained to ensure the patient is fit for general anesthesia and surgery. Routine blood tests (complete blood count, electrolytes, renal and liver function, coagulation profile, blood type and screen) are performed.
  \item \textbf{Medication Adjustments:} Patients are typically instructed to discontinue blood-thinning medications (e.g., aspirin, NSAIDs, warfarin, novel oral anticoagulants) several days to a week prior to surgery, as advised by their surgeon and cardiologist. PPIs are often continued until the day of surgery.
  \item \textbf{NPO Status:} Patients must be NPO (nothing by mouth) for at least 6-8 hours before surgery to minimize the risk of aspiration during anesthesia.
  \item \textbf{Bowel Preparation:} While not always required, some surgeons may recommend a clear liquid diet for 1-2 days prior to surgery, especially for patients with significant constipation.
  \item \textbf{Antibiotic Prophylaxis:} Intravenous broad-spectrum antibiotics (e.g., cefazolin) are administered within 60 minutes prior to skin incision to reduce the risk of surgical site infection.
  \item \textbf{Informed Consent:} Detailed discussion with the surgeon regarding the procedure, potential benefits, risks, alternatives (including continued medical management or partial fundoplication), and expected recovery.
\end{itemize}

\section*{Contraindications \& Precautions}
Careful patient selection is crucial to optimize outcomes and minimize complications.

\textbf{Absolute Contraindications:}
\begin{itemize}
  \item \textbf{Severe Esophageal Dysmotility:} Conditions such as achalasia, severe diffuse esophageal spasm, or advanced scleroderma esophagus, where the esophagus cannot effectively propel food. A 360-degree wrap in these cases is highly likely to cause severe and persistent dysphagia.
  \item \textbf{Uncorrectable Coagulopathy:} Significant bleeding disorders that cannot be managed preoperatively, posing an unacceptable risk of hemorrhage.
  \item \textbf{Inability to Tolerate General Anesthesia:} Severe cardiopulmonary compromise or other medical conditions that make general anesthesia excessively risky.
  \item \textbf{Active Esophageal Malignancy:} Fundoplication is not a treatment for cancer; definitive oncologic management takes precedence.
  \item \textbf{Undiagnosed or Unconfirmed GERD:} Surgery should not be performed without objective evidence of pathological reflux (e.g., positive 24-hour pH study).
\end{itemize}

\textbf{Relative Contraindications and Precautions:}
\begin{itemize}
  \item \textbf{Morbid Obesity (BMI > 35-40 kg/m\textasciicircum2):} While not an absolute contraindication, obesity significantly increases technical difficulty, operative time, risk of complications (e.g., wound infection, DVT), and may lead to higher failure rates of the fundoplication. Bariatric surgery may be considered first or concurrently.
  \item \textbf{Previous Upper Abdominal Surgery:} May result in significant adhesions, increasing operative time and risk of injury to adjacent organs. A laparoscopic approach may be more challenging or require conversion to open.
  \item \textbf{Short Esophagus:} A truly short esophagus, often seen in severe chronic reflux with stricture, may lead to tension on the wrap and potential for wrap migration or disruption. This may necessitate an esophageal lengthening procedure (e.g., Collis gastroplasty) or a partial fundoplication.
  \item \textbf{Severe Cardiopulmonary Disease:} Increases anesthetic and surgical risks, requiring thorough preoperative optimization.
  \item \textbf{Uncontrolled Psychiatric Illness or Chronic Narcotic Use:} May affect patient compliance, pain perception, and satisfaction with surgical outcomes.
  \item \textbf{Bougie Size:} Using an appropriately sized bougie (typically 54-60 Fr) during wrap creation is crucial to prevent the wrap from being too tight, which can cause severe postoperative dysphagia.
  \item \textbf{Vagal Nerve Injury:} Care must be taken during dissection around the esophagus to avoid injury to the vagal nerves, which can lead to gastric emptying issues (gastroparesis) or diarrhea.
  \item \textbf{Splenic Injury:} Mobilization of the gastric fundus and division of short gastric vessels carries a risk of splenic capsular tear or injury, potentially requiring splenectomy, especially in cases of dense adhesions.
\end{itemize}

\section*{Complications \& Risks}
As with any surgical procedure, Nissen fundoplication carries potential risks, which can be broadly categorized into general surgical risks and procedure-specific complications.

\textbf{General Surgical Risks:}
\begin{itemize}
  \item \textbf{Bleeding:} Intraoperative or postoperative hemorrhage, potentially requiring blood transfusion or reoperation.
  \item \textbf{Infection:} Surgical site infection (SSI) at port sites, intra-abdominal abscess, or pneumonia.
  \item \textbf{Anesthesia Complications:} Adverse reactions to medications, respiratory or cardiac problems, aspiration pneumonitis.
  \item \textbf{Deep Vein Thrombosis (DVT) / Pulmonary Embolism (PE):} Blood clots forming in the legs that can travel to the lungs, prevented by prophylactic measures.
  \item \textbf{Injury to Adjacent Organs:} Accidental injury to the liver, spleen, stomach, intestines, or major blood vessels during instrument manipulation or dissection.
\end{itemize}

\textbf{Procedure-Specific Risks:}
\begin{itemize}
  \item \textbf{Dysphagia (Difficulty Swallowing):} The most common complication, often transient due to esophageal edema, but can be persistent if the wrap is too tight or due to underlying esophageal dysmotility. May require endoscopic dilation or, rarely, revision surgery.
  \item \textbf{Gas-Bloat Syndrome:} Inability to belch or vomit due to the competent wrap, leading to abdominal distension, discomfort, and flatulence. Usually resolves over time as the wrap loosens slightly, but can be bothersome.
  \item \textbf{Wrap Herniation or Disruption:} The fundoplication can slip back into the chest (wrap migration) or unravel over time, leading to recurrence of reflux symptoms. This often necessitates revision surgery.
  \item \textbf{Esophageal or Gastric Injury/Perforation:} Perforation of the esophagus or stomach during dissection or suturing, a serious complication requiring immediate repair and potentially prolonged hospitalization.
  \item \textbf{Vagal Nerve Injury:} Can lead to delayed gastric emptying (gastroparesis), chronic diarrhea, or abdominal discomfort.
  \item \textbf{Pneumothorax:} Accidental entry into the pleural cavity during dissection, leading to lung collapse, usually managed with a chest tube insertion.
  \item \textbf{Recurrence of Reflux:} Despite a technically successful operation, some patients may experience a return of GERD symptoms, often due to wrap failure, disruption, or progression of underlying disease.
  \item \textbf{Diarrhea:} Can occur due to vagal nerve irritation or injury, or altered gastric emptying.
  \item \textbf{Abdominal Pain:} Persistent or chronic abdominal pain unrelated to the incision sites, sometimes due to nerve entrapment or adhesions.
\end{itemize}

\section*{Postoperative Recovery Timeline}
Immediately after Nissen fundoplication, patients are transferred to the Post-Anesthesia Care Unit (PACU) for close monitoring of vital signs, pain, and nausea. Pain management is initiated, often with intravenous analgesics, and antiemetics are given to prevent nausea and vomiting. Most patients are able to ambulate within a few hours of surgery, which is encouraged to prevent complications like deep vein thrombosis.

The first 24-48 hours typically involve a gradual advancement of diet. Patients usually start with clear liquids, progressing to a full liquid diet, and then a soft, pureed diet. It is crucial to avoid solid foods initially to prevent stress on the newly formed wrap and minimize dysphagia. Intravenous fluids are continued until oral intake is adequate. Most patients are discharged home within 1-2 days, provided their pain is controlled, they are tolerating liquids, and they are ambulating independently. For the first 1-2 weeks, patients remain on a soft or pureed diet, gradually introducing soft, moist solid foods. Strenuous activity, heavy lifting (over 10-15 lbs), and abdominal exercises should be avoided for 4-6 weeks to allow the surgical sites to heal and prevent disruption of the wrap. Full recovery and resolution of initial postoperative dysphagia or gas-bloat syndrome can take several weeks to months, with most patients returning to normal activities and diet within 6-8 weeks.

\section*{Surgical Findings \& Pathology}
During a Nissen fundoplication, the surgeon documents several key intraoperative findings. These include the presence and size of any hiatal hernia (e.g., Type I sliding hernia, Type II paraesophageal hernia), the condition of the phrenoesophageal ligament, the degree of esophageal inflammation (esophagitis), and the presence of any pre-existing conditions like Barrett's esophagus. The length of the intra-abdominal esophagus achieved after mobilization, the integrity of the diaphragmatic crura, and the successful creation of a tension-free, 360-degree fundic wrap over an appropriate bougie are all critical observations. Any intraoperative complications, such as splenic capsular tears or esophageal perforations, are also meticulously recorded.

If biopsies were taken during preoperative endoscopy (e.g., from areas of esophagitis or suspected Barrett's esophagus) or if a hiatal hernia sac was excised, these tissues are sent for pathological examination. The pathology report will confirm the histological nature of the esophageal and gastric mucosa, detailing the presence and severity of esophagitis (e.g., acute or chronic inflammation, erosions), intestinal metaplasia (Barrett's esophagus), or dysplasia. It will also confirm the benign nature of any resected hiatal hernia sac, typically composed of fibrous and adipose tissue. The pathology report serves to confirm the preoperative diagnosis and rule out any underlying malignancy, providing crucial information for long-term follow-up.

\section{Expected Postoperative Course}
A normal, successful postoperative course following Nissen fundoplication is characterized by a progressive and predictable recovery. In the immediate postoperative period, patients will experience incisional pain, which should be well-controlled with oral analgesics. Transient dysphagia (difficulty swallowing) and gas-bloat syndrome are common but expected to gradually improve over the first few weeks to months. Patients should be able to tolerate a liquid diet on the day of surgery or the next day, progressing to a soft/pureed diet within the first week, and gradually reintroducing solid foods over 4-6 weeks. Nausea and vomiting should be minimal and responsive to antiemetics. Bowel function typically returns within 1-3 days. The surgical incisions should heal cleanly without signs of infection. Ultimately, the patient should experience significant relief from their preoperative GERD symptoms, such as heartburn and regurgitation, and be able to discontinue or significantly reduce their reliance on anti-reflux medications, leading to an improved quality of life.

\section{Postoperative Care \& Follow-up}
Postoperative follow-up is essential to monitor recovery, address any complications, and assess the long-term efficacy of the Nissen fundoplication.
\begin{itemize}
  \item \textbf{Initial Postoperative Visits:} Typically, the first follow-up visit occurs 1-2 weeks after surgery to check wound healing, remove sutures or staples if necessary, and discuss dietary progression and activity restrictions. A subsequent visit may be scheduled at 1-3 months to assess symptom resolution and overall recovery.
  \item \textbf{Dietary Counseling:} Ongoing guidance on diet advancement is crucial, emphasizing small, frequent meals, thorough chewing, and avoiding carbonated beverages or foods that cause discomfort. Management of potential side effects like dysphagia or gas-bloat syndrome is discussed.
  \item \textbf{Activity Resumption:} Patients are advised on a gradual return to normal activities, with specific instructions to avoid heavy lifting (typically >10-15 lbs) or strenuous abdominal exercises for 4-6 weeks to prevent disruption of the wrap.
  \item \textbf{Medication Management:} Patients are typically advised to discontinue PPIs post-surgery, but may be prescribed a short course of antacids or H2 blockers for residual symptoms or during the healing phase.
  \item \textbf{Imaging/Labs:} Routine imaging or laboratory tests are not typically required unless there are specific concerns. If symptoms recur or new symptoms develop, further investigations may include:
    \begin{itemize}
      \item \textbf{Upper Endoscopy:} To evaluate the integrity of the wrap, assess for esophagitis, or rule out other pathologies.
      \item \textbf{Barium Swallow Study:} To visualize the anatomy of the wrap, detect wrap herniation, or identify strictures.
      \item \textbf{Esophageal Manometry and pH Monitoring:} To objectively assess wrap function and confirm recurrent reflux.
    \end{itemize}
  \item \textbf{Warning Signs:} Patients are advised to contact their surgeon immediately if they experience severe or worsening abdominal pain, persistent vomiting, inability to swallow even liquids, fever, chills, redness or pus from incision sites, or any signs of bleeding.
\end{itemize}

\section*{Epidemiology}
Gastroesophageal reflux disease (GERD) is a highly prevalent condition, affecting approximately 10-20\% of the adult population in Western countries. While most cases are managed effectively with lifestyle modifications and medical therapy, a significant subset of patients develops severe or refractory GERD, making them candidates for surgical intervention. Nissen fundoplication is the most commonly performed anti-reflux surgery worldwide. The incidence of fundoplication procedures peaked in the early 2000s following the widespread adoption of the laparoscopic technique, with thousands performed annually. Patients undergoing Nissen fundoplication typically range from young adults to the elderly, with a slight male predominance in those presenting with severe complications like Barrett's esophagus. The most frequent indications for surgery include medically refractory GERD, large hiatal hernias, and significant regurgitation. Mortality rates for laparoscopic Nissen fundoplication are exceedingly low, generally less than 0.1-0.2\% in experienced centers. Morbidity rates, encompassing complications like dysphagia, gas-bloat syndrome, and reoperation for wrap failure, range from 5-15\%.

\section*{Outcomes \& Prognosis}
The long-term outcomes of Nissen fundoplication are generally excellent, with high rates of symptom control and patient satisfaction. Studies report that 85-90\% of patients experience significant improvement or complete resolution of their GERD symptoms at 1 year, and 70-80\% maintain good symptom control at 5-10 years. The ability to discontinue or significantly reduce PPI use is achieved in 80-90\% of patients. Functional outcomes, such as improved quality of life and reduced impact of GERD on daily activities, are also consistently reported. However, a small percentage of patients (5-15\%) may experience recurrence of reflux symptoms over time, often due to wrap failure, disruption, or migration, which may necessitate revision surgery. Persistent dysphagia or gas-bloat syndrome can also affect long-term satisfaction. Prognostic factors for successful outcomes include a positive response to preoperative PPI therapy, objective evidence of reflux on pH monitoring, normal esophageal motility on manometry, and the absence of severe obesity. While Nissen fundoplication effectively controls reflux, it does not typically reverse established Barrett's esophagus, although it may prevent its progression. Overall, for carefully selected patients, Nissen fundoplication offers a durable and effective solution for chronic GERD.

\section*{References}
\begin{enumerate}
  \item MedlinePlus. Nissen fundoplication. U.S. National Library of Medicine; 2024. Available from: https://medlineplus.gov/ency/article/007324.htm
  \item Brunicardi FC, Andersen DK, Billiar TR, et al. Schwartz's Principles of Surgery. 11th ed. McGraw-Hill Education; 2019. Chapter 29: Esophagus and Diaphragm.
  \setcounter{enumi}{2}
  \item Kahrilas PJ, Shaheen RF, Vaezi MM, et al. American Gastroenterological Association Institute Technical Review on the Management of Gastroesophageal Reflux Disease. Gastroenterology. 2008;135(4):1392-1413.
  \setcounter{enumi}{3}
  \item SAGES Guidelines Committee. SAGES Guidelines for the Surgical Treatment of Gastroesophageal Reflux Disease (GERD). Society of American Gastrointestinal and Endoscopic Surgeons; 2021. Available from: https://www.sages.org/publications/guidelines/guidelines-for-the-surgical-treatment-of-gastroesophageal-reflux-disease-gerd/
\end{enumerate}

\clearpage